%% Elevark fra mattelab.no – laget automatisk av bygg.py. Rediger gjerne til eget bruk.
%% Kompiler med lualatex.
\documentclass[11pt,a4paper]{article}
\usepackage[norsk]{babel}
\usepackage[a4paper,margin=2.2cm,top=2cm,bottom=2.4cm]{geometry}
\usepackage{fontspec}
\usepackage{unicode-math}
\setmainfont{TeX Gyre Pagella}
\setsansfont{TeX Gyre Heros}[Scale=MatchLowercase]
\setmathfont{TeX Gyre Pagella Math}
\usepackage{xcolor}
\definecolor{blekk}{HTML}{16213A}
\definecolor{dempet}{HTML}{566077}
\definecolor{rute}{HTML}{CFC9B8}
\definecolor{m1}{HTML}{2F4AB0}
\definecolor{m2}{HTML}{99266D}
\usepackage[most]{tcolorbox}
\usepackage{qrcode}
\usepackage{needspace}
\usepackage{fancyhdr}
\usepackage[hidelinks]{hyperref}
\pdfvariable suppressoptionalinfo \numexpr 1+2+4+8+16+32+64+128+256+512\relax
\setlength{\parindent}{0pt}
\setlength{\parskip}{0.4em}
\pagestyle{fancy}
\fancyhf{}
\renewcommand{\headrulewidth}{0pt}
\fancyfoot[L]{\sffamily\footnotesize\color{dempet}mattelab.no · Algebrabrikker · elevark uten fasit}
\fancyfoot[R]{\sffamily\footnotesize\color{dempet}Side \thepage}

\newcommand{\fase}[2]{\tcbox[on line, colframe=#2, colback=white, coltext=#2, boxrule=0.6pt, arc=2pt,
  boxsep=0pt, left=3pt, right=3pt, top=1.5pt, bottom=1.5pt, fontupper=\sffamily\bfseries\footnotesize]{#1}}
\newcommand{\steg}[4]{%
  \par\noindent{\color{rute}\rule{\linewidth}{0.4pt}}\par\nopagebreak
  \noindent\begin{minipage}[t]{2.1cm}\vspace{0pt}\fase{#1}{#2}\end{minipage}%
  \begin{minipage}[t]{\dimexpr\linewidth-2.1cm\relax}\vspace{0pt}#3\end{minipage}\par
  \vspace{#4}}
\newcommand{\del}[4]{%
  \par\noindent{\color{rute}\rule{\linewidth}{0.4pt}}\par\nopagebreak
  \noindent\begin{minipage}[t]{1.1cm}\vspace{0pt}\fase{#1}{#2}\end{minipage}%
  \begin{minipage}[t]{\dimexpr\linewidth-1.1cm\relax}\vspace{0pt}#3\end{minipage}\par
  \vspace{#4}}
\newcommand{\oppgave}[2]{\par\Needspace{12\baselineskip}\vspace{0.6em}%
  {\fontsize{14}{17}\selectfont\bfseries{\color{m1}#1}\enspace #2}\par\vspace{0.2em}}

\begin{document}
\begin{minipage}[t]{0.72\linewidth}\vspace{0pt}
{\sffamily\small\color{dempet}Elevark · 1T · mattelab.no}\par\vspace{0.3em}
{\fontsize{24}{28}\selectfont\bfseries Algebrabrikker}\par\vspace{1em}
{\sffamily\small Navn: \rule{4.6cm}{0.4pt}\quad Klasse: \rule{1.5cm}{0.4pt}\quad Dato: \rule{1.9cm}{0.4pt}}
\end{minipage}\hfill
\begin{minipage}[t]{0.22\linewidth}\vspace{0pt}\raggedleft
\qrcode[height=2.3cm]{https://mattelab.no/utforskere/algebrabrikker.html}\par\vspace{0.2em}
{\sffamily\scriptsize\color{dempet}Utforskeren}
\end{minipage}

\vspace{1em}
{\sffamily\small\color{dempet}Utforskeren ligger på \url{https://mattelab.no/utforskere/algebrabrikker.html}}\par
Hver oppgave har tre steg. \textbf{Gjett} først, uten å bruke utforskeren. \textbf{Test} så i utforskeren. \textbf{Forklar} til slutt med egne ord hvorfor det blir slik.

\oppgave{1}{En parentes ganget med $x$}
\steg{Gjett}{m1}{Hvilke brikker trenger du for å lage et rektangel med sider $x$ og $x + 3$?}{2.2cm}
\steg{Test}{m2}{Vis $x(x + 3)$. Tell brikkene.}{2.6cm}
\steg{Forklar}{blekk}{Skriv produktet som et uttrykk uten parentes, og forklar det med brikkene.}{3.4cm}

\oppgave{2}{Hvor kommer $5x$ fra?}
\steg{Gjett}{m1}{Gang ut $(x + 2)(x + 3)$ for hånd. Hvor mange $x$-brikker tror du det blir?}{2.2cm}
\steg{Test}{m2}{Vis $(x + 2)(x + 3)$. Finn $x$-brikkene. Hvor ligger de?}{2.6cm}
\steg{Forklar}{blekk}{Forklar hvorfor tallet foran $x$ er $2 + 3$, mens konstantleddet er $2 \cdot 3$.}{3.4cm}

\oppgave{3}{Et kvadrat av brikker}
\steg{Gjett}{m1}{Mange skriver $(x + 3)^2 = x^2 + 9$. Hvilke brikker mangler i så fall?}{2.2cm}
\steg{Test}{m2}{Vis $(x + 3)(x + 3)$. Er figuren et kvadrat? Hvor mange $x$-brikker er det?}{2.6cm}
\steg{Forklar}{blekk}{Sett $x = 1$ i $(x + 3)^2$ og i $x^2 + 9$. Bruk svaret og figuren til å forklare feilen.}{3.4cm}

\oppgave{4}{Nullpar}
\steg{Gjett}{m1}{Hvor mange $x$-brikker blir det i $(x + 3)(x - 3)$? Gjett før du tester.}{2.2cm}
\steg{Test}{m2}{Vis $(x + 3)(x - 3)$. Finn de positive og de negative $x$-brikkene.}{2.6cm}
\steg{Forklar}{blekk}{Forklar hvorfor det ikke blir noe $x$-ledd i svaret. Hvilken setning er dette?}{3.4cm}

\oppgave{5}{Med tall foran $x$ og negative tall}
\steg{Gjett}{m1}{Gang ut $(2x - 1)(x + 4)$ for hånd.}{2.2cm}
\steg{Test}{m2}{Vis $(2x - 1)(x + 4)$. Hvor mange positive og negative $x$-brikker er det? Hvor mange nullpar?}{2.6cm}
\steg{Forklar}{blekk}{Kontroller svaret ditt ved å sette inn $x = 1$ på begge sider.}{3.4cm}

\end{document}
